\section{模型假设}

以下假设只用于补足题面未描述的执行细节；题面已明确的误差界、距离阈值和计时规则作为模型约束保留。每项同时说明采用依据及其对模型的影响。
\begin{itemize}
  \item \textbf{运动环境：}题面未设置障碍物或定位误差，故假设机器人按已接受指令到达目标坐标。由此可用欧氏距离核算移动时间；若实际环境存在障碍，需要将欧氏距离替换为可行路径距离，并重新检查检测点的可达性，定位观测模型本身保持不变。
  \item \textbf{观测误差：}示向误差仅满足题设 $[-1^\circ,1^\circ]$ 有界条件，且同一地点误差固定。不额外假设独立同分布，也不以重复测量平均降噪，因而所有定位结论均按最坏情形成立。
  \item \textbf{源状态：}一次任务内，源的位置、频道、接收半径、类型和发射方向保持不变，每个频道至多对应一个源。该假设使历史观测可以持续求交；若源会移动，则需改为动态状态估计。
  \item \textbf{接收半径：}将 $R_c\in[1000,1500]\,\mathrm{m}$ 始终作为未知量，不用上界替代真实值。因此问题二的检测点必须满足保证接收约束，全向源负观测也只能排除半径 $1000\,\mathrm{m}$ 的圆盘。
  \item \textbf{动作一致性：}仅在服务器确认接受后更新位置、频道和虚拟时间；网络重试复用同一请求标识。该约定防止通信重试被误计为重复动作。
  \item \textbf{边界处理：}理论推导采用题面给出的闭边界，浮点实现加入不改变物理阈值含义的微小保护量，并用边界样例复核，以避免数值舍入造成假阳性证书。
\end{itemize}
